\chapter{Mecánica Relativista}

\section{Tensores de Minkowski}

El tensor $A^\mu$, de orden $(1,0)$, de acuerdo con la ley de
transformación de tensores bajo una transformación de coordenadas
$x'^\mu=x'^\mu(x^\nu)$, cambia como
\[
A'^\mu = \frac{\partial x'^\mu}{\partial x^\nu}\,A^\nu .
\]
En el espacio de Minkowski $M^4$, las transformaciones de coordenadas
son las transformaciones de Lorentz, $x'^\mu=\Lambda^\mu_{\ \nu}(\beta^k)
x^\nu$, por lo tanto
\[
\frac{\partial x'^\mu}{\partial x^\nu} = \Lambda^\mu_{\ \nu}(\beta^k),
\tag{1}
\]
tal que el tensor $A^\mu$ en $M^4$ transforma como
\[
A'^\mu = \Lambda^\mu_{\ \nu}(\beta^k)\,A^\nu . \tag{2}
\]
Usando la definición de $\Lambda^\mu_{\ \nu}(\beta^k)$ en (2) se obtiene
\[
A'^0 = \gamma\big(A^0+\beta_iA^i\big), \qquad A'^i = \gamma\big(\beta^i
A^0+\hat n_jA^j\hat n^i\big)+\epsilon^{iln}\hat n_l\epsilon_{njk}A^j
\hat n^k,
\]
o, en forma vectorial,
\[
A'^0 = \gamma\big(A^0+\vec\beta\cdot\vec A\big), \qquad \vec A\,' =
\gamma\big(\vec\beta A^0+(\hat n\cdot\vec A)\hat n\big)+\hat n\times
(\vec A\times\hat n). \tag{3}
\]
Las componentes de $A^\mu$ paralelas a $\hat n$ son las afectadas por las
transformaciones de Lorentz (dilataciones y contracciones), mientras que
las perpendiculares permanecen invariantes.

Para el tensor $A_\mu$, de orden $(0,1)$, que cambia como
$A'_\mu=(\partial x^\nu/\partial x'^\mu)A_\nu$, se tiene que, bajo
transformaciones de Lorentz, $x'^\mu=\Lambda^\mu_{\ \nu}(\beta^k)x^\nu
\Rightarrow x^\mu=(\Lambda^{-1})^\mu_{\ \nu}x'^\nu$, luego
\[
\frac{\partial x^\mu}{\partial x'^\nu} = (\Lambda^{-1})^\mu_{\ \nu}
(\beta^k) = \Lambda^\mu_{\ \nu}(-\beta^k), \tag{4}
\]
tal que el tensor $A_\mu$ en $M^4$ transforma
\[
A'_\mu = (\Lambda^{-1})^\nu_{\ \mu}(\beta^k)A_\nu = \Lambda^\nu_{\ \mu}
(-\beta^k)A_\nu . \tag{5}
\]
Usando la definición de $\Lambda^\nu_{\ \mu}(-\beta^k)$ en (5),
\[
A'_0 = \gamma\big(A_0-\beta^iA_i\big), \qquad A'_i = \gamma\big(-\beta_i
A_0+\hat n^jA_j\hat n_i\big)+\epsilon_{iln}\hat n^l\epsilon^{njk}A_j
\hat n_k,
\]
y en forma vectorial
\[
A'_0 = \gamma\big(A_0-\vec\beta\cdot\vec A\big), \qquad \vec A\,' =
\gamma\big(-\vec\beta A_0+(\hat n\cdot\vec A)\hat n\big)+\hat n\times
(\vec A\times\hat n). \tag{6}
\]
Nuevamente, las componentes paralelas a $\hat n$ son las afectadas por
la transformación de Lorentz y las perpendiculares permanecen
invariantes.

En general, un tensor $T^{\mu_1\dots\mu_p}_{\quad\ \nu_1\dots\nu_q}$ de
orden $(p,q)$, bajo una transformación general de coordenadas
$x'^\mu=x'^\mu(x^\nu)$, cambia como
\[
T'^{\mu_1\dots\mu_p}_{\quad\ \nu_1\dots\nu_q} = \frac{\partial
x'^{\mu_1}}{\partial x^{\alpha_1}}\cdots\frac{\partial
x'^{\mu_p}}{\partial x^{\alpha_p}}\,\frac{\partial x^{\beta_1}}{\partial
x'^{\nu_1}}\cdots\frac{\partial x^{\beta_q}}{\partial x'^{\nu_q}}\,
T^{\alpha_1\dots\alpha_p}_{\quad\ \beta_1\dots\beta_q}.
\]
Para el caso particular del espacio de Minkowski, las transformaciones de
coordenadas deben satisfacer (1) y (4). Por lo tanto, un tensor de orden
$(p,q)$ en el espacio de Minkowski debe transformar con $p$
transformaciones de Lorentz y $q$ transformaciones inversas de Lorentz.

\begin{definicion}[Tensor de Minkowski]
Un tensor $T^{\mu_1\dots\mu_p}_{\quad\ \nu_1\dots\nu_q}$ de orden
$(p,q)$ se dice un \textbf{tensor de Minkowski} si transforma como
\[
T'^{\mu_1\dots\mu_p}_{\quad\ \nu_1\dots\nu_q} = \Lambda^{\mu_1}_{\
\alpha_1}\cdots\Lambda^{\mu_p}_{\ \alpha_p}\,(\Lambda^{-1})^{\beta_1}_{\
\nu_1}\cdots(\Lambda^{-1})^{\beta_q}_{\ \nu_q}\,T^{\alpha_1\dots
\alpha_p}_{\quad\ \beta_1\dots\beta_q}. \tag{7}
\]
\end{definicion}

Con el fin de pasar de la mecánica clásica (no relativista) a la
mecánica relativista debemos exigir que las magnitudes físicas sean
tensores de Lorentz (incluyendo escalares). Por ejemplo, la velocidad
$\vec v$ es una cantidad de la mecánica clásica que debe generalizarse a
un espacio de cuatro dimensiones: la nueva velocidad tendrá cuatro
componentes, se llamará \textbf{cuadrivelocidad} y debe ser un tensor de
Lorentz. La segunda ley de Newton también debe generalizarse,
introduciendo una \textbf{cuadrifuerza} con cuatro componentes, y la
segunda ley de Newton relativista debe escribirse de forma covariante de
Lorentz.

\section{Sistema de referencia propio}

\subsection{El tiempo propio}

Sea una partícula $P$ que se mueve respecto de un S.R.I. $K$, tal que su
posición queda determinada por el cuadrivector posición $x^\mu$, con
$\mu=0,1,2,3$. La partícula, al moverse en el espaciotiempo
cuadridimensional, genera una línea de mundo. Un cambio infinitesimal de
posición consiste en pasar de $x^\mu$ a $x^\mu+dx^\mu$, de modo que el
cambio de posición es el tensor $dx^\mu$, que es tangente a la línea de
mundo.

\begin{center}
\begin{tikzpicture}[scale=1]
  \draw[-Stealth] (0,0) -- (4.4,0) node[right] {$x^1$};
  \draw[-Stealth] (0,0) -- (0,4.4) node[above] {$x^0$};
  \draw[thick,colordef] (1.6,0.2) .. controls (2.3,1.4) and (1.4,2.6) .. (2.2,4.0);
  \node[colordef,right] at (2.3,3.6) {\footnotesize línea de mundo};
  \fill (1.8625,2.025) circle (1.5pt) node[left] {$P$};
  \draw[-Stealth,thick,colorprop] (1.8625,2.025) -- (1.81,2.85) node[right] {$dx^\mu$};
\end{tikzpicture}
\end{center}

Para definir una cuadrivelocidad necesitamos una variable temporal que
permita definir un antes y un después, tal que
\[
u^\mu = \frac{\text{cambio en } x^\mu}{\text{cambio en el tiempo}},
\]
y la definición debe ser tal que $u^\mu$ sea un tensor de Lorentz de
orden $(1,0)$, esto es, debe transformar bajo una transformación de
Lorentz como
\[
u'^\mu = \Lambda^\mu_{\ \nu}\,u^\nu . \tag{8}
\]
La tentación inicial es usar el tiempo medido en el S.R.I. $K$, es decir
definir
\[
u^\mu = \frac{dx^\mu}{dt} . \tag{9}
\]
Sin embargo, el tiempo $t$ de $K$ es relativo: si cambiamos al S.R.I.
$K'$,
\[
dx'^\mu = \Lambda^\mu_{\ \nu}dx^\nu, \qquad dt' = \frac1c\Lambda^0_{\
\nu}dx^\nu = \Lambda^0_{\ 0}dt + \frac1c\Lambda^0_{\ j}dx^j,
\]
luego
\[
u'^\mu = \frac{dx'^\mu}{dt'} = \frac{\Lambda^\mu_{\ \nu}dx^\nu}{
\Lambda^0_{\ 0}dt+\frac1c\Lambda^0_{\ j}dx^j} = \Lambda^\mu_{\ \nu}
\frac{u^\nu}{\Lambda^0_{\ 0}+\Lambda^0_{\ j}\frac{v^j}{c}} \neq
\Lambda^\mu_{\ \nu}u^\nu .
\]
Por lo tanto, la definición (9) no satisface la condición de covariancia
de Lorentz (8). Para salvar la definición de la cuadrivelocidad como un
tensor de Lorentz introduciremos el llamado \textbf{sistema de
referencia propio}.

Sea una partícula $P$ que se mueve con velocidad $v^i=dx^i/dt$ respecto
del S.R.I. $K$. Definimos un sistema de referencia atado a la partícula,
$K_P$, no necesariamente inercial: un observador atado a la partícula,
con unidades de medida, es decir, con un reloj, reglas y balanza. Para
entender la idea, recordemos una de las aventuras del barón Münchausen,
en la que viaja sobre una bala de cañón: la bala sería la partícula bajo
estudio y el barón el observador propio que viaja junto con la
partícula.

\begin{definicion}[Sistema coordenado propio]
Se define un \textbf{sistema coordenado propio}, denotado $K_P$, como un
observador que se mueve junto con la partícula, tal que la partícula
está en reposo respecto de $K_P$, y hay un reloj asociado a $K_P$ en
reposo respecto de la partícula que mide lo que se llama el
\textbf{tiempo propio}, denotado $t_P=\tau$.
\end{definicion}

Las coordenadas de la partícula respecto de $K_P$ son
\[
x^\mu_P = \big(x^0_P,x^i_P\big) = (c\tau,\vec 0), \qquad x^i_P=0, \tag{10}
\]
y la condición $x^i_P=0$ es válida $\forall\tau$, pues el sistema $K_P$
se mueve junto a la partícula, es decir, la partícula está en el origen
de $K_P$. El elemento de línea respecto de $K_P$ es
\[
ds_P^2 = \eta_{\mu\nu}dx^\mu_Pdx^\nu_P = -(dx_P^0)^2+dx^i_Pdx^i_P =
-c^2d\tau^2<0,
\]
es decir, es de tipo tiempo.

Después de que la partícula se ha movido durante un intervalo de tiempo
medido por el reloj del S.R.I. $K$, ha recorrido una distancia espacial
$dl=\sqrt{dx^2+dy^2+dz^2}$ medida en $K$, y el elemento de línea en $K$
es
\[
ds^2 = -c^2dt^2+dl^2 = -c^2dt^2+dx^2+dy^2+dz^2. \tag{11}
\]
En $K_P$ el reloj marca un intervalo $d\tau$ y la distancia espacial
recorrida por la partícula es $dl_P=\sqrt{dx_P^2+dy_P^2+dz_P^2}=0$, por
lo que
\[
ds_P^2 = -c^2d\tau^2 . \tag{12}
\]

\begin{proposicionc}[Tiempo propio]
El tiempo medido en $K$ y el tiempo propio están relacionados por
\[
d\tau = dt\sqrt{1-\frac{v^2}{c^2}} . \tag{13}
\]
\end{proposicionc}

\begin{proof}
Los elementos de línea (11) y (12) deben ser iguales, $ds_P^2=ds^2$:
\[
-c^2d\tau^2 = -c^2dt^2+dx^2+dy^2+dz^2 \;\Rightarrow\;
c^2d\tau^2 = c^2dt^2-dx^idx^i = c^2dt^2\left(1-\frac{v^iv_i}{c^2}\right),
\]
de donde $d\tau=dt\sqrt{1-v^2/c^2}$. \qedhere
\end{proof}

La expresión (13) relaciona el tiempo medido en $K$ con el tiempo
propio, si se conoce la velocidad $v^i$ medida en $K$ (que puede variar
con el tiempo, $v^i(t)$). Además, como $ds_P^2=ds^2$ implica
$d\tau=\frac1c\sqrt{-ds^2}$, y el elemento de línea medido en $K$ debe
ser de tipo tiempo ($ds^2<0$),
\[
d\tau = \frac{ds}{c}, \qquad ds:=\sqrt{-ds^2}. \tag{14}
\]

\begin{proposicionc}
El tiempo propio es invariante bajo transformaciones de Lorentz: es un
escalar de Lorentz.
\end{proposicionc}

\begin{proof}
Como el elemento de línea es invariante bajo transformaciones de
Lorentz, si efectuamos sucesivas transformaciones de Lorentz $K
\xrightarrow{V^i}K'\xrightarrow{V'^i}K''\to\cdots$ se tiene
$d\tau=\frac1cds=\frac1cds'=\frac1cds''=\dots$, es decir
$d\tau=d\tau'=d\tau''=\dots$ \qedhere
\end{proof}

El tiempo propio es el tiempo que se mide en el sistema propio: si una
mosca tiene una vida media de 28 días en su sistema propio, no importa
si la mosca vuela a altas velocidades (una ``supermosca'') respecto de
otros S.R.I.; su vida media en su sistema propio siempre será de 28
días, aunque los observadores inerciales observen dilatación del tiempo.
Como el tiempo propio es un invariante de Lorentz, lo podemos usar para
medir la secuencia de eventos en el espaciotiempo $M^4$: la línea de
mundo se parametriza con el tiempo propio de la partícula,
\[
\mathbb{R}\to M^4, \qquad \tau\mapsto x^\mu(\tau),
\]
y el cambio infinitesimal de posición se expresa como
$dx^\mu=x^\mu(\tau+d\tau)-x^\mu(\tau)=\dfrac{dx^\mu}{d\tau}d\tau$.

\subsection{La cuadrivelocidad}

\begin{definicion}[Cuadrivelocidad]
Se define la \textbf{cuadrivelocidad} de una partícula como
\[
u^\mu := \frac{dx^\mu}{d\tau} . \tag{15}
\]
\end{definicion}

\begin{proposicionc}
La cuadrivelocidad (15) es un tensor de Lorentz de orden $(1,0)$.
\end{proposicionc}

\begin{proof}
Sea la transformación de Lorentz $x'^\mu=\Lambda^\mu_{\ \nu}x^\nu$;
entonces $dx'^\mu=\Lambda^\mu_{\ \nu}dx^\nu$ y $d\tau'=d\tau$, de modo
que
\[
u'^\mu = \frac{dx'^\mu}{d\tau'} = \Lambda^\mu_{\ \nu}\frac{dx^\nu}{d\tau}
= \Lambda^\mu_{\ \nu}u^\nu,
\]
es decir, la cuadrivelocidad satisface la propiedad (8). \qedhere
\end{proof}

La cuadrivelocidad es tangente a la línea de mundo de la partícula. Para
partículas físicas debe estar dentro del cono de luz, si $v<c$, o sobre
el cono de luz, si $v=c$. Se debe mencionar que la cuadrivelocidad
combina una cantidad medida por el S.R.I. $K$, $dx^\mu$, y una cantidad
medida por el sistema propio $K_P$, $d\tau$:
\[
u^\mu = \frac{\text{cambio de posición medido en } K}{\text{cambio en el
tiempo medido en } K_P}.
\]
El tiempo propio no es una cantidad medida directamente por $K$, pero
ambas están relacionadas a través de (13). Si conocemos $v(t)$,
integrando (13),
\[
\tau(t) = \tau_0+\int_{t_0}^t\sqrt{1-\frac{v^2(t)}{c^2}}\,dt,
\]
y si la partícula se mueve con velocidad constante respecto de $K$,
$v=V$, se tiene $\tau=\sqrt{1-V^2/c^2}\,t$, es decir, $t=\gamma\tau$.

Conocida la velocidad $v^i=dx^i/dt$ medida por $K$, podemos
relacionarla con la cuadrivelocidad:
\[
u^\mu = \left(c\frac{dt}{d\tau},\frac{dx^i}{d\tau}\right) =
\left(c\frac{dt}{d\tau},\frac{dt}{d\tau}\frac{dx^i}{dt}\right),
\]
y usando (13),
\[
u^0 = \frac{c}{\sqrt{1-v^2/c^2}} = c\gamma, \qquad u^i =
\frac{v^i}{\sqrt{1-v^2/c^2}} = \gamma v^i, \qquad
\gamma=\frac{1}{\sqrt{1-v^2/c^2}}. \tag{16, 17}
\]

\begin{proposicionc}[Ligadura de la cuadrivelocidad]
La cuadrivelocidad satisface la ligadura
\[
u^2 = \eta_{\mu\nu}u^\mu u^\nu = -c^2, \tag{18}
\]
o, en otra notación, $u^\mu u_\mu=-c^2$ (19).
\end{proposicionc}

\begin{proof}
Partiendo de la identidad $ds^2=ds_P^2$, $\eta_{\mu\nu}dx^\mu dx^\nu=
-c^2d\tau^2$; como $x^\mu=x^\mu(\tau)$,
$\eta_{\mu\nu}\dfrac{dx^\mu}{d\tau}\dfrac{dx^\nu}{d\tau}d\tau^2=
-c^2d\tau^2$, de donde $\eta_{\mu\nu}u^\mu u^\nu=-c^2$. \qedhere
\end{proof}

La expresión (18) muestra que la cuadrivelocidad es de tipo tiempo, y
que no todas sus componentes son independientes: las cuatro componentes
de $u^\mu$ deben tomar solo valores que satisfagan la ligadura (18).
Para entenderlo, consideremos una partícula que se mueve a lo largo del
eje $x$ de $K$, es decir $u^2=u^3=0$; entonces
\[
-(u^0)^2+(u^1)^2=-c^2 \;\Longrightarrow\; (u^0)^2-(u^1)^2=c^2, \tag{20}
\]
que corresponde a una hipérbola en el espacio de las cuadrivelocidades.

\begin{center}
\begin{tikzpicture}[scale=1]
  \draw[-Stealth] (-3,0) -- (3,0) node[right] {$u^1$};
  \draw[-Stealth] (0,-3.2) -- (0,3.2) node[above] {$u^0$};
  \draw[thick,colordef,domain=-2.2:2.2,samples=60,smooth]
    plot ({\x},{1.0*sqrt(1+\x*\x)});
  \draw[thick,colordef,domain=-2.2:2.2,samples=60,smooth]
    plot ({\x},{-1.0*sqrt(1+\x*\x)});
  \draw[dashed,gray] (-2.4,-2.4) -- (2.4,2.4);
  \draw[dashed,gray] (-2.4,2.4) -- (2.4,-2.4);
  \fill (1.2,1.562) circle (1.5pt) node[right] {$P$};
  \fill (0.6,0.5) circle (1.5pt) node[right] {$Q$};
  \node[left] at (0,1.0) {\footnotesize$c$};
\end{tikzpicture}
\end{center}

Las componentes $u^\mu=(u^0,u^1,0,0)$ deben tomar solo valores sobre la
hipérbola (la rama positiva o la negativa). La cuadrivelocidad puede
tomar los valores del punto $P$ sobre la hipérbola, pero no los del
punto $Q$, que no está sobre ella. De (20), $u^1=\sqrt{(u^0)^2-c^2}$, y
es claro que $|u^0|\ge c>0$: la componente $u^0>0$ nunca se anula. Si una
persona se sienta en una silla está en reposo, $u^1=0$, por lo tanto
$u^0=c$: nos movemos en la dirección temporal, envejeciendo.

\subsection{La cuadrivelocidad reducida}

\begin{definicion}[Cuadrivelocidad reducida]
Como $ds$ es de tipo tiempo, $ds^2<0$, podemos usarlo para definir un
nuevo cuadrivector, la \textbf{cuadrivelocidad reducida},
\[
U^\mu := \frac{dx^\mu}{ds}, \tag{21}
\]
que es adimensional.
\end{definicion}

\begin{proposicionc}
La cuadrivelocidad reducida es un tensor de Lorentz de orden $(1,0)$.
\end{proposicionc}

\begin{proof}
Sea $x'^\mu=\Lambda^\mu_{\ \nu}x^\nu$, tal que $dx'^\mu=\Lambda^\mu_{\
\nu}dx^\nu$ y $ds'=ds$. Entonces $U'^\mu=dx'^\mu/ds'=\Lambda^\mu_{\
\nu}dx^\nu/ds=\Lambda^\mu_{\ \nu}U^\nu$, que satisface la propiedad (2).
\qedhere
\end{proof}

Usando (14) en (21), $U^\mu=\dfrac{dx^\mu}{ds}=\dfrac1c\dfrac{dx^\mu}
{d\tau}$, de modo que
\[
U^\mu = \frac{u^\mu}{c}, \tag{22}
\]
expresión que relaciona la cuadrivelocidad con la reducida. Las
unidades son distintas: $[u^\mu]=L/T$, $[U^\mu]=1$. Calculemos $U^2$:
\[
U^2 = \eta_{\mu\nu}U^\mu U^\nu = \frac{1}{c^2}\eta_{\mu\nu}u^\mu u^\nu =
-\frac{c^2}{c^2} = -1,
\]
usando (18). La cuadrivelocidad reducida es de tipo tiempo; desde el
punto de vista geométrico, $U^\mu$ es el cuadrivector unitario tangente
a la línea de mundo.

\subsection{La cuadriaceleración}

\begin{definicion}[Cuadriaceleración]
Si $u^\mu$ es la cuadrivelocidad, se define la \textbf{cuadriaceleración}
como
\[
W^\mu := \frac{du^\mu}{d\tau} = \frac{d^2x^\mu}{d\tau^2} . \tag{23}
\]
\end{definicion}

\begin{proposicionc}
La cuadriaceleración es un tensor de Lorentz de orden $(1,0)$.
\end{proposicionc}

\begin{proof}
Sea $x'^\mu=\Lambda^\mu_{\ \nu}x^\nu$, tal que $u'^\mu=\Lambda^\mu_{\
\nu}u^\nu$ y $d\tau'=d\tau$. Entonces $W'^\mu=\dfrac{du'^\mu}{d\tau'}=
\dfrac{d}{d\tau}\big(\Lambda^\mu_{\ \nu}u^\nu\big)=\Lambda^\mu_{\
\nu}\dfrac{du^\nu}{d\tau}=\Lambda^\mu_{\ \nu}W^\nu$. \qedhere
\end{proof}

Al igual que en la cuadrivelocidad, en la definición de la
cuadriaceleración el cambio de tiempo corresponde al tiempo propio y no
al tiempo del S.R.I. $K$. Sin embargo, podemos relacionar $W^\mu$ con el
tiempo de $K$: $W^\mu=\dfrac{du^\mu}{d\tau}=\dfrac{dt}{d\tau}\dfrac{
du^\mu}{dt}=\gamma\dfrac{du^\mu}{dt}$. Usando (16),
\[
W^0 = \gamma\frac{d}{dt}(c\gamma) = c\gamma\frac{d\gamma}{dt} =
\frac c2\frac{d\gamma^2}{dt} = \frac c2\frac{d}{dt}\left(1-\frac{v^2}
{c^2}\right)^{-1} = \frac1c\left(1-\frac{v^2}{c^2}\right)^{-2}v_i
\frac{dv^i}{dt},
\]
\[
W^i = \gamma\frac{d}{dt}\big(\gamma v^i\big) = \gamma^2\frac{dv^i}{dt}+
\gamma v^i\frac{d\gamma}{dt} = \gamma^2\frac{dv^i}{dt} + \frac{1}{c^2}
\left(1-\frac{v^2}{c^2}\right)^{-2}v^iv_j\frac{dv^j}{dt}.
\]

\begin{proposicionc}[Cuadriaceleración en términos de $\vec v$ y $\vec a$]
Con $v^i=dx^i/dt$ y $a^i=dv^i/dt$,
\[
W^0 = \frac1c\gamma^4\,\vec v\cdot\vec a, \qquad W^i = \gamma^2a^i +
\frac{1}{c^2}\gamma^4(\vec v\cdot\vec a)\,v^i . \tag{24}
\]
\end{proposicionc}

Se ha relacionado $W^\mu$ con las magnitudes físicas medidas por $K$.

\begin{proposicionc}
La cuadrivelocidad y la cuadriaceleración son perpendiculares:
\[
u\cdot W = \eta_{\mu\nu}u^\mu W^\nu = 0, \tag{25}
\]
o, en otra notación, $u_\mu W^\mu=0$ (26).
\end{proposicionc}

\begin{proof}
Derivamos la ligadura (18): $\dfrac{d}{d\tau}\big(\eta_{\mu\nu}u^\mu
u^\nu\big)=-\dfrac{dc^2}{d\tau}=0$, es decir $2\eta_{\mu\nu}u^\mu
\dfrac{du^\nu}{d\tau}=0$, o bien $u\cdot W=0$. \qedhere
\end{proof}

La expresión (25) es una ligadura para la cuadriaceleración: no todas
las componentes de $W^\mu$ son independientes. En efecto,
\[
\eta_{\mu\nu}u^\mu W^\nu=0 \;\Rightarrow\; -u^0W^0+u^iW^i=0
\;\Rightarrow\; -\gamma cW^0+\gamma v_iW^i=0,
\]
\[
\therefore\quad W^0 = \frac1c\,v_iW^i = \frac1c\,\vec v\cdot\vec W .
\tag{27}
\]
Si se reemplazan las expresiones (24) en (27) se obtiene una identidad.

\section{Momentum lineal en el espacio de Minkowski}

En mecánica de Newton el momentum lineal se define como
\[
p^i = mv^i, \tag{28}
\]
donde $m$ es la masa inercial, que mide la oposición del sistema a
cambiar su estado de movimiento.

\begin{definicion}[Cuadrimomentum lineal]
Por analogía con (28) se define el \textbf{cuadrimomentum lineal} como
\[
p^\mu := m_0u^\mu, \tag{29}
\]
donde $m_0$ es la masa medida en el sistema propio $K_P$.
\end{definicion}

La masa $m_0$ medida en $K_P$ es la masa que se mide en reposo, y por eso
se llama \textbf{masa en reposo}. Es una magnitud independiente del
observador, es decir, un invariante bajo cambio de sistema de
referencia, y por eso también se llama \textbf{masa invariante}: si una
persona tiene una masa en reposo de $90\ \text{kg}$, no importa la
velocidad con que se mueva, su masa propia no cambia.

\begin{proposicionc}
El cuadrimomentum lineal (29) es un tensor de Lorentz de orden $(1,0)$.
\end{proposicionc}

\begin{proof}
Sea $x'^\mu=\Lambda^\mu_{\ \nu}x^\nu$, tal que $u'^\mu=\Lambda^\mu_{\
\nu}u^\nu$ y $m_0'=m_0$. Entonces $p'^\mu=m_0'u'^\mu=m_0\Lambda^\mu_{\
\nu}u^\nu=\Lambda^\mu_{\ \nu}p^\nu$. \qedhere
\end{proof}

Cambiando del tiempo propio al tiempo del S.R.I. $K$,
$p^\mu=m_0u^\mu=m_0(c\gamma,\gamma v^i)$, así
\[
p^0 = m_0c\gamma, \qquad p^i = m_0\gamma v^i . \tag{30}
\]
Si comparamos (28) con la segunda ecuación de (30) podemos definir la
nueva magnitud
\[
m = m(v) = m_0\gamma = \frac{m_0}{\sqrt{1-v^2/c^2}} . \tag{31}
\]
La masa (31) es la masa medida en el S.R.I. $K$, y claramente no es
constante; se llama \textbf{masa relativista}.

\begin{center}
\begin{tikzpicture}[scale=1]
  \draw[-Stealth] (0,0) -- (6.4,0) node[right] {$v$};
  \draw[-Stealth] (0,0) -- (0,4.6) node[above] {$m(v)$};
  \draw[dashed] (5,0) -- (5,4.4);
  \node[below] at (5,0) {\footnotesize$c$};
  \node[left] at (0,0.8) {\footnotesize$m_0$};
  \draw[thick,colordef,domain=0:0.98,samples=80,smooth]
    plot ({5*\x},{0.8/sqrt(1-\x*\x)});
\end{tikzpicture}
\end{center}

En mecánica clásica se introduce la masa inercial, que es una constante.
En relatividad especial también hay una masa inercial, que también mide
la oposición a cambiar el estado de movimiento, pero varía con la
velocidad: aumenta a medida que aumenta la velocidad. La razón es la
constancia de la rapidez de la luz y, como consecuencia, la existencia
de una velocidad máxima: a medida que aumenta la velocidad la masa
inercial también aumenta, teóricamente hasta el infinito si $v\to c$,
para que ninguna partícula supere la velocidad de la luz. La masa
inercial es, por así decirlo, la ``policía'' de la relatividad especial:
evita que la partícula viole el segundo postulado de Einstein.

\begin{observacion}
Consideremos una persona de masa en reposo $90\ \text{kg}$ que subimos a
un vehículo espacial, aumentando la velocidad suavemente hasta alcanzar
velocidades cercanas a la de la luz. El observador $K$ observará que la
masa relativista de la persona aumenta de acuerdo con (31). Sin embargo,
esto no significa que $K$ vea que la persona se vuelve obesa. Lo que
observa es que la persona se mueve cada vez más lento a medida que
aumenta la velocidad, debido a la dilatación del tiempo; el observador
interpreta ese efecto de moverse lento como un aumento de la inercia, y
como la inercia se relaciona con la masa, interpreta este incremento
como un aumento de la masa relativista.
\end{observacion}

Las componentes del cuadrimomentum lineal no son todas independientes,
pues satisfacen la ligadura
\[
p^2 = \eta_{\mu\nu}p^\mu p^\nu = m_0^2\eta_{\mu\nu}u^\mu u^\nu = -m_0^2c^2
. \tag{32}
\]
El cuadrimomentum es un cuadrivector de tipo tiempo. La expresión (32)
se denomina \textbf{capa de masa} (en inglés, \emph{mass shell}): dada
una masa $m_0$, las componentes del cuadrimomentum tomarán los valores
tales que (32) se satisfaga. Esto es particularmente importante en teoría
cuántica de campos. Desarrollando (32),
\[
-(p^0)^2+p^ip^i = -m_0^2c^2 \;\Longrightarrow\; (p^0)^2-\|\vec p\|^2 =
m_0^2c^2,
\]
que es un hiperboloide de dos hojas en el espacio de los
cuadrimomentum (con $p^0>0$ solo la hoja superior); además se concluye
que $p^0\neq0$.

\section{Segunda ley de Newton relativista}

Las ecuaciones de Newton son invariantes de forma (covariantes) bajo
transformaciones de Galileo, pero no bajo transformaciones de Lorentz.
En mecánica relativista deben generalizarse para satisfacer los
postulados de Einstein, y deben desembocar en la segunda ley de Newton
para velocidades muy pequeñas comparadas con la de la luz. Esto se
resume en dos puntos:
\begin{enumerate}
\item Las ecuaciones de Newton generalizadas deben ser tensores de
Lorentz (covariancia de Lorentz).
\item Para $v\ll c$ deben desembocar en la segunda ley de Newton,
\[
v\ll c \;\Longrightarrow\; \frac{d}{dt}\big(m_0v^i\big) = F^i, \tag{33}
\]
porque la teoría de Newton ha sido exitosa en explicar la naturaleza a
pequeñas velocidades.
\end{enumerate}

\begin{definicion}[Ecuaciones de Newton relativistas]
Las ecuaciones de Newton relativistas son de la forma
\[
\frac{dp^\mu}{d\tau} = K^\mu, \tag{34}
\]
donde $K^\mu$ es la \textbf{cuadrifuerza}, también llamada \textbf{fuerza
de Minkowski}.
\end{definicion}

\begin{proposicionc}
La fuerza de Minkowski $K^\mu$ es un tensor de Lorentz de orden $(1,0)$.
\end{proposicionc}

\begin{proof}
Sea $x'^\mu=\Lambda^\mu_{\ \nu}x^\nu$, tal que $p'^\mu=\Lambda^\mu_{\
\nu}p^\nu$ y $d\tau'=d\tau$. Introduciendo en (34),
$K'^\mu=\dfrac{dp'^\mu}{d\tau'}=\dfrac{d}{d\tau}\big(\Lambda^\mu_{\
\nu}p^\nu\big)=\Lambda^\mu_{\ \nu}\dfrac{dp^\nu}{d\tau}=\Lambda^\mu_{\
\nu}K^\nu$. \qedhere
\end{proof}

Para identificar las componentes de $K^\mu$ con las observaciones
realizadas por el S.R.I. $K$ debemos pasar del tiempo propio al tiempo
del observador $K$:
\[
K^\mu = \frac{dp^\mu}{d\tau} = \frac{dt}{d\tau}\frac{dp^\mu}{dt} =
\gamma\frac{dp^\mu}{dt} \;\Rightarrow\; \big(K^0,K^i\big) = \left(\gamma
\frac{d}{dt}(m_0c\gamma),\ \gamma\frac{d}{dt}(m_0\gamma v^i)\right),
\tag{35}
\]
e identificando los términos de cada lado,
\[
K^0 = \gamma\frac{dp^0}{dt} = \gamma\frac{d}{dt}(m_0c\gamma), \qquad
K^i = \gamma\frac{dp^i}{dt} = \gamma\frac{d}{dt}\big(m_0\gamma v^i\big).
\tag{36}
\]
Comparando (28) con la segunda ecuación de (36) realizamos la
identificación
\[
\frac{dp^i}{dt} = \frac{d}{dt}\big(m_0\gamma v^i\big) = \frac{d}{dt}
\big(m(v)v^i\big) = F^i . \tag{37}
\]
Esta identificación se debe a dos razones. La primera es que, para
$v\ll c$, se debe desembocar en la segunda ley de Newton no
relativista. En efecto, desarrollando $\gamma$ en serie,
\[
\gamma = \frac{1}{\sqrt{1-v^2/c^2}} = 1+\frac{v^2}{2c^2}+\frac{3v^4}
{8c^4}+\dots
\]
y entonces
\[
F^i = \frac{d}{dt}\big(m_0\gamma v^i\big) = \frac{d}{dt}\big(m_0v^i\big)
+ \frac{d}{dt}\left(\frac{m_0v^iv^2}{2c^2}+\frac{3m_0v^iv^4}{8c^4}+\dots
\right) = \frac{d}{dt}\big(m_0v^i\big)+\dots,
\]
de modo que, para velocidades pequeñas comparadas con la de la luz,
$F^i\approx\dfrac{d}{dt}(m_0v^i)$, que concuerda con (33). La segunda
razón es conceptual: en mecánica clásica el momentum contiene a la masa
inercial $m$, que es constante porque no hay límite superior a la
velocidad. En relatividad hay una velocidad máxima, y por lo tanto la
inercia crece a medida que aumenta la velocidad, y este crecimiento se
debe a la prohibición de superar la velocidad de la luz; como la masa es
la medida de la inercia, la masa inercial debe también crecer:
\[
p^i = \frac{d}{dt}(m_0v^i) \;\xrightarrow{\text{relativista}}\; p^i =
\frac{d}{dt}\left(\frac{m_0}{\sqrt{1-v^2/c^2}}\,v^i\right).
\]
Introduciendo (37) en la segunda ecuación de (36),
\[
K^i = \gamma F^i . \tag{38}
\]
Desarrollando la primera ecuación de (36),
\[
K^0 = \gamma\frac{d}{dt}(m_0c\gamma) = cm_0\gamma\frac{d\gamma}{dt} =
\frac{m_0c}{2}\frac{d\gamma^2}{dt} \;\Rightarrow\; K^0 = \frac{m_0}{c}
\gamma^4\,\vec v\cdot\vec a . \tag{39}
\]
Desarrollando la segunda ecuación de (36),
\[
K^i = \gamma\frac{d}{dt}\big(m_0\gamma v^i\big) = m_0\gamma^2\frac{dv^i}
{dt}+\frac{m_0}{2}v^i\frac{d\gamma^2}{dt} = \Big(m_0\gamma^2\delta^i_j+
\frac{m_0}{c^2}\gamma^4v^iv_j\Big)a^j,
\]
es decir,
\[
K^i = \gamma\Big(m_0\gamma\delta^i_j+\frac{m_0}{c^2}\gamma^3v^iv_j\Big)
a^j . \tag{40}
\]
Definimos la cantidad de masa (\textbf{no es un tensor})
\[
M^i_{\ j} = m_0\gamma\,\delta^i_j+\frac{m_0}{c^2}\gamma^3v^iv_j,
\tag{41}
\]
de modo que reemplazando (41) en (40),
\[
K^i = \gamma M^i_{\ j}a^j . \tag{42}
\]
Para obtener información de (41) escribamos la velocidad como
$v^i=v\hat n^i$ (43), donde $\hat n^i$ es el vector unitario en la
dirección de la velocidad y $v$ la rapidez. Entonces
\[
M^i_{\ j} = m_0\gamma\delta^i_j+\frac{m_0v^2}{c^2}\gamma^3\hat n^i\hat
n_j = m_0\gamma\big(\delta^i_j-\hat n^i\hat n_j\big)+\Big(\frac{m_0v^2}
{c^2}\gamma^3+m_0\gamma\Big)\hat n^i\hat n_j,
\]
y como $\dfrac{v^2}{c^2}+\gamma^{-2}=1$, se obtiene
\[
M^i_{\ j} = m_0\gamma\big(\delta^i_j-\hat n^i\hat n_j\big)+m_0\gamma^3
\hat n^i\hat n_j = m_0\gamma P^i_{\ j}+m_0\gamma^3\hat n^i\hat n_j .
\tag{44, 45}
\]
El primer término es proporcional al operador proyector $P^i_{\ j}$, por
lo que es perpendicular a la velocidad y posee una masa $m_0\gamma$; el
segundo es paralelo a la velocidad y posee una masa $m_0\gamma^3$. Para
entender su significado físico pasemos al S.R.I. $K$: introduciendo (38)
en (42), $\gamma F^i=\gamma M^i_{\ j}a^j$, es decir
\[
F^i = M^i_{\ j}a^j . \tag{46}
\]
La aceleración puede descomponerse en una parte paralela y otra
perpendicular a la velocidad, $a^i=a^i_\parallel+a^i_\perp$. Como
$a^i=\dfrac{d}{dt}(v\hat n^i)$,
\[
a^i = \frac{dv}{dt}\hat n^i+v\frac{d\hat n^i}{dt} . \tag{47}
\]
El término $\hat n^i\,dv/dt$ es paralelo a la velocidad. El término
$v\,d\hat n^i/dt$ es perpendicular a ella; en efecto, contrayendo con
$\hat n_i$,
\[
\hat n_i\frac{d\hat n^i}{dt} = \frac12\frac{d}{dt}(\hat n^2) = \frac12
\frac{d}{dt}(1) = 0. \tag{48}
\]
Resumiendo,
\[
a^i_\parallel = \frac{dv}{dt}\hat n^i, \qquad a^i_\perp = v\frac{d\hat
n^i}{dt}. \tag{49, 50}
\]
Reemplazando en (46), $F^i=M^i_{\ j}a^j_\parallel+M^i_{\ j}a^j_\perp$.
Calculamos cada término:
\[
M^i_{\ j}a^j_\parallel = \big(m_0\gamma P^i_{\ j}\hat n^j+m_0\gamma^3
\hat n^i\hat n^2\big)\frac{dv}{dt} = m_0\gamma^3a^i_\parallel, \tag{51}
\]
\[
M^i_{\ j}a^j_\perp = m_0\gamma\big(\delta^i_j-\hat n^i\hat n_j\big)
\frac{d\hat n^j}{dt}v = m_0\gamma\,a^i_\perp, \tag{52}
\]
usando (48). Reemplazando (51) y (52) en (46),
\[
F^i = m_0\gamma^3a^i_\parallel+m_0\gamma\,a^i_\perp . \tag{53}
\]
En mecánica clásica, la razón entre fuerza y aceleración es la masa
inercial. La expresión (53) no cumple con esto: hay una masa en la
dirección paralela y una masa diferente en la dirección perpendicular a
la velocidad. En la dirección paralela aparece la \textbf{masa
longitudinal}
\[
m_l = m_0\gamma^3 = \frac{m_0}{(1-v^2/c^2)^{3/2}}, \tag{54}
\]
y en la dirección perpendicular aparece la \textbf{masa transversal}
\[
m_t = m_0\gamma = \frac{m_0}{\sqrt{1-v^2/c^2}} . \tag{55}
\]
Como $\gamma^3\ge\gamma$, se tiene $m_l\ge m_t$: cuesta más acelerar la
partícula en la dirección paralela a la velocidad que en la
perpendicular.

\begin{proposicionc}
La cuadrifuerza y la cuadrivelocidad son perpendiculares:
\[
u\cdot K = u_\mu K^\mu = 0 . \tag{56}
\]
\end{proposicionc}

\begin{proof}
\[
u_\mu K^\mu = \eta_{\mu\nu}u^\nu\frac{dp^\mu}{d\tau} = m_0\eta_{\mu\nu}
u^\nu\frac{du^\mu}{d\tau} = \frac{m_0}{2}\frac{d}{d\tau}\big(\eta_{\mu\nu}
u^\mu u^\nu\big) = -\frac{m_0}{2}\frac{dc^2}{d\tau} = 0 . \qedhere
\]
\end{proof}

La ecuación (56) es una ligadura para la cuadrifuerza: no todas las
componentes $K^\mu$ son independientes. Desarrollándola,
$-u^0K^0+u^iK^i=0$, es decir $-c\gamma K^0+\gamma v_iK^i=0$, y por lo
tanto
\[
K^0 = \frac1c v_iK^i = \frac1c\,\vec v\cdot\vec K . \tag{57}
\]
Introduciendo (38) en (57) se obtiene $K^0=\dfrac{\gamma}{c}F^iv_i=
\dfrac\gamma c\vec F\cdot\vec v$. La cantidad $\vec F\cdot\vec v$ es la
potencia $\mathcal P$, luego
\[
K^0 = \frac\gamma c\mathcal P = \frac\gamma c\frac{dW}{dt}, \tag{58}
\]
donde $W$ es el trabajo mecánico. La componente $K^0$ está dada también
por la primera ecuación de (36); igualando,
\[
\frac\gamma c\frac{dW}{dt} = \gamma\frac{d}{dt}(m_0c\gamma)
\;\Longrightarrow\; \frac{dW}{dt} = \frac{d}{dt}\big(m_0c^2\gamma\big).
\tag{59}
\]
Como la fuerza es la fuerza neta, por el teorema de las fuerzas vivas el
trabajo mecánico es igual a la energía cinética $T$. Luego
\[
\frac{dT}{dt} = \frac{d}{dt}\big(m_0c^2\gamma\big) = \frac{d}{dt}\left(
\frac{m_0c^2}{\sqrt{1-v^2/c^2}}\right). \tag{60}
\]
Integrando (60) con respecto del tiempo,
\[
T = \frac{m_0c^2}{\sqrt{1-v^2(t)/c^2}}-\frac{m_0c^2}{\sqrt{1-v^2(0)/c^2}},
\tag{61}
\]
y para $t=0$ se escoge $v(0)=0$, de modo que
\[
T = \frac{m_0c^2}{\sqrt{1-v^2(t)/c^2}}-m_0c^2 . \tag{62}
\]
Debemos notar que hemos calculado la energía cinética de una partícula
de masa en reposo $m_0$ y masa relativista $m_0\gamma$. Esta expresión
puede escribirse
\[
T+m_0c^2 = m_0c^2\gamma = mc^2 . \tag{63}
\]
El término $m_0c^2$ apareció de la condición de que para $t=0$ la
velocidad de la partícula es cero; por esta razón se llama \textbf{energía
en reposo}. En la integración que permitió obtener (61) se integró desde
$t=0$, $v(0)=0$, pero se pudo integrar desde $t=t_1$ hasta $t=t_2>t_1$,
\[
T = \frac{m_0c^2}{\sqrt{1-v^2(t_2)/c^2}}-\frac{m_0c^2}{\sqrt{1-v^2(t_1)/
c^2}}. \tag{64}
\]
Decimos que la partícula en $t=t_1$ tiene velocidad $v^i_1$ o masa
relativista $m_1=m_0/\sqrt{1-v_1^2/c^2}$, y en $t=t_2$ aumenta su
velocidad a $v^i_2$ o su masa relativista a $m_2$, tal que la energía
cinética es
\[
T = m_2c^2-m_1c^2 = (m_2-m_1)c^2 = \Delta m\,c^2 . \tag{65}
\]
Este es un resultado de gran importancia en la física: relaciona la
energía cinética con la variación de la masa inercial de la partícula a
causa del aumento de la velocidad. Volviendo al caso $v(0)=0$, si la
energía cinética está dada por (63), entonces la energía $E$ de la
partícula libre es $E=T+m_0c^2=mc^2$, es decir,
\[
E = mc^2 = \frac{m_0c^2}{\sqrt{1-v^2/c^2}} . \tag{66}
\]
Este resultado es uno de los más conocidos de la relatividad especial y
establece que \textbf{la masa es una forma de energía}. La energía (66)
se llama \textbf{energía total}: para una partícula libre, está
compuesta de la energía en reposo $m_0c^2$ y la energía cinética $T$.

Le hemos llamado \emph{energía cinética} a la expresión (62); para ver
por qué, consideremos velocidades pequeñas, $v\ll c$, y expandamos:
\[
T = m_0c^2\left(1+\frac12\frac{v^2}{c^2}+\frac38\frac{v^4}{c^4}+\dots
\right)-m_0c^2 = \frac12m_0v^2+\frac38m_0\frac{v^4}{c^2}+\dots \tag{67}
\]
De manera que para velocidades bajas o cotidianas se tiene la
aproximación $T\approx\frac12m_0v^2$, que corresponde a la ecuación no
relativista de la energía cinética. Para una partícula sometida a una
fuerza conservativa, la energía total incluye una energía potencial. La
interpretación de la energía en reposo $m_0c^2$ es la de un nuevo tipo
de energía. La equivalencia entre masa y energía, ecuación (66), ha sido
confirmada experimentalmente en una variedad de casos, por ejemplo en la
fusión nuclear de los núcleos atómicos.

\section{Energía en el espacio de Minkowski}

Retomemos la expresión (57) y desarrollémosla un poco:
\[
K^0 = \frac1cK^iv_i \;\Rightarrow\; \gamma\frac{d}{dt}(m_0c\gamma) =
\frac\gamma cF^iv_i \;\Rightarrow\; \frac{d}{dt}\big(m_0c^2\gamma\big) =
\vec F\cdot\vec v . \tag{68}
\]
En el caso en que la fuerza $F^i$ provenga de un potencial, es decir,
sea una fuerza conservativa, $\vec F=-\vec\nabla U$, entonces
\[
\frac{d}{dt}\big(m_0c^2\gamma\big) = -\vec\nabla U\cdot\vec v = -
\frac{dU}{dt} \;\Longrightarrow\; \frac{d}{dt}\big(m_0c^2\gamma+U\big) =
0. \tag{69}
\]
La integración de (69) conduce a la \textbf{ley relativista de la
energía},
\[
E = m_0c^2\gamma+U = \text{constante}. \tag{70}
\]
La expresión (70) es la energía total de la partícula en presencia de una
fuerza conservativa. Para velocidades bajas o cotidianas se aproxima a
$E\approx m_0c^2+\frac12m_0v^2+U+\dots$ (la constante $m_0c^2$ puede
absorberse en la definición de la energía).

Volvamos al caso de la partícula libre, $U=0$. La componente temporal del
cuadrimomentum es
\[
p^0 = m_0u^0 = m_0\frac{dx^0}{d\tau} = m_0c\gamma = mc = \frac{mc^2}{c} =
\frac Ec, \qquad\text{es decir,}\qquad p^0 = \frac Ec , \tag{71}
\]
donde se usó (66): la componente temporal del cuadrimomentum representa
la energía total de la partícula libre. Volvamos a la componente espacial
de la segunda ley de Newton relativista, $K^i=\dfrac{dp^i}{d\tau}$, es
decir $\gamma F^i=\dfrac{d}{d\tau}\Big(\dfrac{m_0v^i}{\sqrt{1-v^2/c^2}}
\Big)$. Si la trifuerza que actúa sobre la partícula es nula, $F^i=0$,
entonces
\[
\frac{m_0v^i}{\sqrt{1-v^2/c^2}} = \text{constante}, \tag{72}
\]
que es la forma relativista de la ley de conservación del momentum lineal
para una partícula libre. Esta ecuación vectorial implica que la
dirección de $v^i$ es constante. Además, multiplicando (72) por sí misma,
$\dfrac{m_0^2v^2}{1-v^2/c^2}=\text{constante}^2$, de modo que la rapidez
es constante: por lo tanto, la velocidad es constante.

Sea ahora la componente temporal de la segunda ley relativista,
$K^0=\dfrac{dp^0}{d\tau}$, con $p^0=E/c$: $\dfrac{d}{d\tau}(m_0c\gamma)=
\dfrac{d}{d\tau}\Big(\dfrac Ec\Big)$, es decir
\[
\frac{d}{d\tau}\big(m_0c^2\gamma\big) = \frac{dE}{d\tau}\quad\Longrightarrow
\quad E=m_0c^2\gamma=\text{constante}, \tag{73}
\]
que es la ley de conservación de la energía para la partícula libre. Si
hay una fuerza conservativa, la conservación de la energía está dada por
(70).

Volviendo a la partícula libre, escribamos el cuadrimomentum de modo que
en la componente temporal aparezca la energía total,
\[
p^\mu = \left(\frac Ec,p^i\right), \tag{74}
\]
y calculemos $p^2$:
\[
p^2 = \eta_{\mu\nu}p^\mu p^\nu = -m_0^2c^2 = -\frac{E^2}{c^2}+\|\vec
p\|^2 \;\Longrightarrow\; E^2 = m_0^2c^4+c^2\|\vec p\|^2, \quad
E=\sqrt{m_0^2c^4+c^2\|\vec p\|^2}. \tag{75}
\]
La expresión (75) relaciona la energía con el cuadrimomentum de una
partícula libre. Aparentemente tenemos dos fórmulas para la energía de
una partícula libre, (66) y (75), pero son equivalentes:
\[
E^2 = m_0^2c^4+c^2m_0\gamma v^im_0\gamma v_i = m_0^2c^4+m_0^2c^2\gamma^2
v^2 = m_0^2c^4\gamma^2\Big(\gamma^{-2}+\frac{v^2}{c^2}\Big) =
m_0^2c^4\gamma^2 \;\Rightarrow\; E=m_0c^2\gamma=mc^2 .
\]
De (66) despejamos
\[
\gamma = \frac{E}{m_0c^2}, \tag{76}
\]
y reemplazando en la segunda ecuación de (30), $p^i=m_0\gamma v^i=
m_0\dfrac{E}{m_0c^2}v^i$, es decir
\[
p^i = \frac{E}{c^2}v^i , \tag{77}
\]
lo que nos dice que la energía tiene momentum. La magnitud del
trimomentum es
\[
\|\vec p\| = \frac{E}{c^2}v . \tag{78}
\]

\begin{proposicionc}
Si una partícula tiene masa en reposo nula, entonces se mueve a la
velocidad de la luz.
\end{proposicionc}

\begin{proof}
Por hipótesis $m_0=0$; reemplazando en (75),
\[
E = c\|\vec p\| . \tag{79}
\]
Introduciendo (79) en (78), $\|\vec p\|=\dfrac{c\|\vec p\|}{c^2}v$, de
donde $v=c$. \qedhere
\end{proof}

La mecánica cuántica nos enseña que una partícula se comporta como onda
y como partícula. Einstein, para explicar el efecto fotoeléctrico,
postuló que la luz se comporta como una partícula, y nosotros aquí
consideramos la luz como una partícula de masa en reposo igual a cero, a
la que llamaremos \textbf{fotón}. El cuadrimomentum lineal del fotón lo
designaremos $k^\mu$, en vez de $p^\mu$, tal que
\[
k^2 = \eta_{\mu\nu}k^\mu k^\nu = -m_0^2c^2 = 0 , \tag{80}
\]
es decir, el cuadrimomentum del fotón es de tipo luz. Debe escribirse en
la forma (74). Esta forma se dedujo a partir de la partícula libre; sin
embargo, es más general e incluye, por ejemplo, el caso de partículas de
masa en reposo nula.

En la obtención del cuadrimomentum primero se definió la cuadrivelocidad
(15), donde el tiempo propio cumple un papel preponderante. Sin embargo,
para una partícula de tipo luz se cumple que
\[
ds^2=0 \;\Rightarrow\; ds=0, \tag{81}
\]
y entonces, por (14), $d\tau=0$: no se puede definir una cuadrivelocidad
a partir del tiempo propio, pues la expresión (15) se indetermina.
Definir la cuadrivelocidad de una partícula de tipo luz es complicado,
porque hay que buscar un nuevo invariante asociado a ella que permita
definir un antes y un después. Es por ello que el cuadrimomentum se
vuelve fundamental, y asumiremos que la definición (74) es válida para
cualquier tipo de partícula física.

La energía del fotón está relacionada con su frecuencia en la forma
\[
E = h\nu, \tag{82}
\]
donde
\[
h \approx 6.626\times10^{-34}\ \frac{\text{kg}\,\text{m}^2}{\text{s}}
\tag{83}
\]
es la constante de Planck. El cuadrimomentum del fotón se escribe
\[
k^\mu = \left(\frac Ec,k^i\right) = \left(\frac{h\nu}{c},k^i\right).
\tag{84}
\]
Introduciendo (84) en (80): $-\dfrac{h^2\nu^2}{c^2}+k^ik^i=0$, luego
\[
\|\vec k\| = \frac{h\nu}{c} = \frac h\lambda, \tag{85}
\]
donde $\lambda$ es la longitud de onda. El trimomentum $k^i$ se puede
escribir en la forma
\[
k^i = \frac{h\nu}{c}\hat k^i, \tag{86}
\]
con $\hat k^i$ un vector unitario en la dirección del movimiento del
fotón. Por lo tanto,
\[
k^\mu = \left(\frac{h\nu}{c},\frac{h\nu}{c}\hat k^i\right). \tag{87}
\]
Resumimos: el cuadrimomentum de una partícula con masa en reposo
$m_0\neq0$ está dado por (74) y la energía por (70) si hay presente una
fuerza conservativa; para el fotón, el cuadrimomentum está dado por (87)
y la energía por (82).

Si hay $N$ partículas, la segunda ley de Newton relativista del sistema
se escribe
\[
\frac{dp^\mu_1}{d\tau}+\frac{dp^\mu_2}{d\tau}+\dots+\frac{dp^\mu_N}
{d\tau} = K^\mu, \tag{88}
\]
donde $K^\mu$ es la cuadrifuerza externa sobre el sistema (las fuerzas
internas se anulan por la tercera ley de Newton). Definimos el
\textbf{cuadrimomentum total}
\[
p^\mu = p^\mu_1+p^\mu_2+\dots+p^\mu_N = \sum_{i=1}^Np^\mu_i, \tag{89}
\]
y entonces
\[
\frac{dp^\mu}{d\tau} = K^\mu . \tag{90}
\]
Un caso de interés físico es cuando la cuadrifuerza se anula, $K^\mu=0$:
\[
\frac{dp^\mu}{d\tau}=0 \;\Rightarrow\; p^\mu=\text{constante}, \tag{91}
\]
que también se escribe
\[
p^\mu_{\text{inicial}} = p^\mu_{\text{final}} . \tag{92}
\]
La expresión (92) es análoga al caso de la mecánica clásica para el
estudio de colisiones. Sin embargo, aquí no solo sirve para estudiar
colisiones, sino también, por ejemplo, desintegraciones atómicas,
decaimiento de partículas, etc.

\begin{observacion}[Antipartículas]
En la obtención de la energía (75) se extrajo una raíz cuadrada y se
asumió el signo positivo para la energía; el signo menos se eliminó. Sin
embargo, el físico Paul Dirac, al estudiar campos de espín $1/2$
(fermiones), encontró que la energía negativa aparecía en sus ecuaciones
y no logró eliminarla; por lo tanto decidió estudiarlas y denominó
\textbf{antipartículas} a las partículas con energía negativa. La
predicción de Dirac de la existencia de antimateria fue confirmada
experimentalmente por Carl D.~Anderson, en Caltech, quien descubrió el
positrón en 1932. Veintitrés años después, en 1955, Emilio Segrè y Owen
Chamberlain, en la Universidad de Berkeley, descubrieron el antiprotón y
el antineutrón. En 1933 se otorgó el premio Nobel de física a Paul
Dirac.
\end{observacion}

\begin{ejemplo}[Una partícula cargada con velocidad constante no emite un fotón]
Supongamos una partícula cargada que se mueve con velocidad constante y
emite un fotón de frecuencia $\nu$. Sea $p^\mu_i$ su cuadrimomentum
inicial y $p^\mu_f$ el final; las líneas continuas representan
partículas masivas y la línea ondulada un fotón. Como la partícula se
mueve con velocidad constante, usamos un S.R.I. $K$ tal que la partícula
se encuentra en reposo respecto de $K$; entonces su cuadrimomentum
inicial se reduce a
\[
p^\mu_i = \left(\frac{E_i}{c},0\right) = \left(\frac{m_0c^2}{c},0\right)
= (m_0c,0). \tag{93}
\]
Después de emitir el fotón, la energía y la velocidad de la partícula
cambian, $p^\mu_f=(E_f/c,p^i_f)$ (94), y el cuadrimomentum del fotón está
dado por (87). Como la cuadrifuerza externa es cero, de (92),
\[
p^\mu_i = p^\mu_f+k^\mu \quad\Longrightarrow\quad p^\mu_i-k^\mu =
p^\mu_f . \tag{95, 96}
\]
Elevando al cuadrado, y usando $p_i^2=p_f^2=-m_0^2c^2$ y $k^2=0$,
\[
p_i^2-2p_i\cdot k+k^2=p_f^2 \;\Rightarrow\; -m_0^2c^2-2\eta_{\mu\nu}
p^\mu_ik^\nu = -m_0^2c^2 \;\Rightarrow\; \eta_{\mu\nu}p^\mu_ik^\nu=0 .
\]
Desarrollando, $-p^0_ik^0+p^i_ik^i=0$, y con (93),
$-m_0c\,\dfrac{h\nu}{c}+0\cdot k^i=0$, es decir $m_0h\nu=0$, de donde
\[
\nu = 0 .
\]
La partícula no emite un fotón.
\end{ejemplo}

\section{Efecto Compton}

El \textbf{efecto Compton} (o dispersión Compton) consiste en el
aumento de la longitud de onda de un fotón cuando choca con un electrón
libre y pierde parte de su energía. La frecuencia o la longitud de onda
de la radiación dispersada depende únicamente del ángulo de dispersión.
Fue estudiado por el físico Arthur Compton en 1923, quien pudo
explicarlo utilizando la noción cuántica de la radiación
electromagnética como cuantos de energía y la mecánica relativista.
Constituyó la demostración final de la naturaleza cuántica de la luz,
tras los estudios de Planck sobre el cuerpo negro y la explicación de
Einstein del efecto fotoeléctrico.

Compton descubrió este efecto al experimentar con rayos X dirigidos
contra una de las caras de un bloque de carbón. Al chocar, los rayos X se
difundieron en varias direcciones; a medida que el ángulo de los rayos
difundidos aumentaba, también se incrementaba su longitud de onda. Con
base en la teoría cuántica, Compton afirmó que el efecto se debía a que
el cuanto de rayos X actúa como una partícula material al chocar contra
el electrón, por lo cual la energía cinética que el cuanto comunica al
electrón representa una pérdida en su energía original. Como
consecuencia de estos estudios, Compton ganó el Premio Nobel de Física
en 1927.

Este efecto es de especial relevancia científica, ya que no puede ser
explicado a través de la naturaleza ondulatoria de la luz: esta debe
comportarse como partícula para poder explicar dichas observaciones, por
lo que adquiere una dualidad onda--corpúsculo característica de la
mecánica cuántica.

Consideremos un electrón libre. Esto parece un poco raro, si el electrón
está confinado a ser parte de un material (originalmente Compton usó
carbón); sin embargo, los electrones más externos del material se
comportan, en buena aproximación, como si estuvieran libres.

Sea un electrón libre de masa en reposo $m_0=m_e=9.10938291(40)\times
10^{-31}\ \text{kg}$. Escogemos un S.R.I. $K$ en que el electrón está en
reposo, entonces
\[
p^\mu_i = \left(\frac{E_i}{c},p^i_i\right) = \left(\frac{m_ec^2}{c},0
\right) = (m_ec,0).
\]
Un fotón, con frecuencia inicial $\nu_i=\nu$, incide sobre el electrón;
su cuadrimomentum inicial es $k^\mu_i=\Big(\dfrac{h\nu}{c},\dfrac{h\nu}{c}
\hat n^i_i\Big)$, donde $\hat n_i^i$ es la dirección inicial hacia donde
se mueve el fotón.

\begin{center}
\begin{tikzpicture}[scale=1]
  \fill (0,0) circle (2.5pt) node[below left] {$e^-$};
  \draw[-Stealth,thick,colordef] (-3.2,0) -- (-0.15,0)
    node[midway,above] {\footnotesize$k_i^\mu\ (\nu)$};
  \draw[-Stealth,thick,colordef] (0.15,0.08) -- (2.8,1.7)
    node[above right] {\footnotesize$k_f^\mu\ (\nu')$};
  \draw[-Stealth,thick,colorprop] (0.15,-0.08) -- (2.6,-1.5)
    node[below right] {\footnotesize$p_f^\mu$};
  \draw (1.1,0) arc (0:31:1.1);
  \node at (1.45,0.3) {\footnotesize$\theta$};
\end{tikzpicture}
\end{center}

Una vez que el fotón y el electrón interactúan, ambos se dispersan. El
electrón adquiere velocidad, por lo tanto su energía aumenta, y su
cuadrimomentum final es $p^\mu_f=(E_f/c,p^i_f)$. El fotón también cambia
su energía, es decir su frecuencia $\nu_f=\nu'$, y su dirección, tal que
$k^\mu_f=\Big(\dfrac{h\nu'}{c},\dfrac{h\nu'}{c}\hat n^i_f\Big)$, donde
$\hat n^i_f$ es la nueva dirección de propagación del fotón.

Como es una colisión (las fuerzas internas se anulan) y no hay fuerzas
externas, se conserva el cuadrimomentum del sistema, ecuación (92):
\[
p^\mu_f+k^\mu_f = p^\mu_i+k^\mu_i . \tag{97}
\]
La componente temporal de (97), $p^0_f+k^0_f=p^0_i+k^0_i$, da
\[
\frac{E_f}{c}+\frac{h\nu'}{c} = m_ec+\frac{h\nu}{c} \;\Rightarrow\;
E_f+h\nu' = m_ec^2+h\nu , \tag{98}
\]
que es la conservación de la energía. Reescribiendo (97) en la forma
\[
p^\mu_f-p^\mu_i = k^\mu_i-k^\mu_f \tag{99}
\]
y elevando al cuadrado,
\[
p_f^2-2p_f\cdot p_i+p_i^2 = k_f^2-2k_f\cdot k_i+k_i^2 \;\Rightarrow\;
-m_e^2c^2-2p_f\cdot p_i-m_e^2c^2 = -2k_f\cdot k_i,
\]
es decir
\[
p_f\cdot p_i+m_e^2c^2 = k_f\cdot k_i . \tag{100}
\]
Calculemos cada producto. Como $p_i^\mu=(m_ec,0)$,
\[
p_f\cdot p_i = -m_ec\frac{E_f}{c}+0\cdot p^i_f = -m_eE_f, \tag{101}
\]
y
\[
k_f\cdot k_i = -\frac{h^2\nu\nu'}{c^2}+\frac{h^2\nu\nu'}{c^2}\hat n_i\cdot
\hat n_f = -\frac{h^2\nu\nu'}{c^2}\,(1-\cos\theta), \tag{102}
\]
pues $\hat n_i\cdot\hat n_f=\cos\theta$. Introduciendo (101) y (102) en
(100),
\[
m_eE_f-m_e^2c^2 = \frac{h^2\nu\nu'}{c^2}(1-\cos\theta). \tag{103}
\]
De (98) despejamos
\[
E_f = m_ec^2+h(\nu-\nu'), \tag{104}
\]
e introduciéndola en (103),
\[
m_e\big(m_ec^2+h(\nu-\nu')\big)-m_e^2c^2 = \frac{h^2\nu\nu'}{c^2}
(1-\cos\theta) \;\Rightarrow\; m_eh(\nu-\nu') = \frac{h^2\nu\nu'}{c^2}
(1-\cos\theta),
\]
\[
\nu-\nu' = \frac{h\nu\nu'}{m_ec^2}(1-\cos\theta).
\]
Multiplicando por $c/\nu\nu'$ se obtiene $\dfrac c{\nu'}-\dfrac c\nu=
\dfrac h{m_ec}(1-\cos\theta)$, y como $\lambda=c/\nu$,
\[
\lambda'-\lambda = \frac{h}{m_ec}\,(1-\cos\theta). \tag{105}
\]
La ecuación (105) se conoce como \textbf{efecto Compton}.

\begin{definicion}[Longitud de onda de Compton]
Se define la \textbf{longitud de onda de Compton} de una partícula de
masa en reposo $m_0$ como
\[
\lambda_c = \frac{h}{m_0c} . \tag{106}
\]
\end{definicion}

Con esta definición el efecto Compton se reescribe como
\[
\Delta\lambda = \lambda_c\,(1-\cos\theta). \tag{107}
\]
Esta expresión proviene del análisis de la interacción como si fuera una
colisión elástica, y su deducción requiere únicamente la utilización de
los principios de conservación de la energía y del momentum.

\section{Caja de Einstein}

En el experimento pensado por A.~Einstein en 1906, llamado la
\textbf{caja de Einstein}, se estudia la relación entre la inercia de la
materia y la energía de radiación. Estudiaremos qué cantidad de masa
relativista equivale a una energía dada. Para ello suponemos que en el
extremo izquierdo de una caja de masa $M$ y longitud $L$, inicialmente en
reposo, se emite un fotón de energía $E$.

\begin{center}
\begin{tikzpicture}[scale=1]
  \draw[thick] (0,1.6) rectangle (5,2.6);
  \draw[-Stealth,thick,colordef] (0.2,2.1) -- (4.8,2.1)
    node[midway,above] {\footnotesize fotón $E$};
  \node at (2.5,2.85) {\footnotesize antes: caja en reposo};
  \draw[thick] (-0.6,-0.1) rectangle (4.4,0.9);
  \draw[Stealth-Stealth] (-0.6,-0.4) -- (0,-0.4) node[midway,below] {\footnotesize$\Delta x$};
  \node at (1.9,1.15) {\footnotesize después: la caja retrocedió};
\end{tikzpicture}
\end{center}

El fotón lleva un momentum $k=E/c$. Debido a que el momentum total del
sistema caja--fotón debe conservarse, y como inicialmente la caja está
en reposo y no hay fotón, el momentum inicial total es nulo y, por lo
tanto, el momentum total final también se anula; así, con la emisión del
fotón la caja recibe una transferencia de impulso $p=-E/c$. Debido a este
retroceso, la caja se mueve con la velocidad
\[
v = -\frac{E}{Mc} . \tag{108}
\]
Después de un tiempo $\Delta t$, el fotón incide en la pared opuesta de la
caja, la cual se detiene, debido a que el momentum transferido para
frenarla es igual al negativo del impulso inicial. Por tanto, la caja se
desplaza una distancia
\[
\Delta x = v\Delta t = -\frac{E\Delta t}{Mc} . \tag{109}
\]
Por otro lado, el fotón demora el mismo tiempo en recorrer el largo de
la caja,
\[
\Delta t = \frac Lc , \tag{110}
\]
e introduciendo (110) en (109),
\[
\Delta x = -\frac{EL}{Mc^2} . \tag{111}
\]
Si ponemos el centro de masa $x_{CM}$ del sistema en el origen del
sistema coordenado, su posición debe permanecer inalterada después de
terminar el experimento. Entonces $X_{CM,\text{inicial}}=X_{CM,
\text{final}}$, es decir $0=\Delta xM/M$, de donde $\Delta x=0$, lo cual
es una contradicción, pues la caja se movió. Esta contradicción solo es
posible salvarla si atribuimos una masa $m$ al fotón, tal que
\[
\frac{\Delta xM+mL}{M+m} = 0 \;\Longrightarrow\; mL = -\Delta xM .
\tag{112}
\]
Introduciendo (111) en (112), $mL=\dfrac{ELM}{Mc^2}$, es decir
\[
E = mc^2 . \tag{113}
\]
La ecuación (113) describe la inercia de la energía: cualquier cambio $E$
de la energía de un cuerpo provoca un cambio correspondiente $\Delta m$
de su masa relativista. En el extremo de la caja donde se emite el fotón,
la masa relativista se reduce en $E/c^2$. En consecuencia, la masa
relativista de la caja aumenta de nuevo en la misma cantidad cuando el
fotón se detiene o se termaliza en el otro extremo de la caja.
